% \subsubsection{Ângulos de Euler}
% \label{sec:angulos_euler}

% Assim como a matriz de cossenos diretores, a orientação de um referencial $B$
% em relação a outro $N$ pode ser descrita por ângulos de Euler.
% Conforme \cite{wie2008space}, considera-se a sequência de rotações
% $C_1(\phi_1) \leftarrow C_2(\phi_2) \leftarrow C_3(\phi_3)$ que leva o
% referencial inercial $N$ até o sistema do corpo $B$. Introduzem-se dois
% referenciais intermediários $N'$ e $N''$:

% \begin{equation}
% C_3(\phi_3): N' \leftarrow N,\qquad
% C_2(\phi_2): N'' \leftarrow N',\qquad
% C_1(\phi_1): B \leftarrow N'',
% \end{equation}

% Isto é, parte-se do referencial inercial $N$ e aplica-se primeiro $C_3(\phi_3)$, obtendo $N'$, depois $C_2(\phi_2)$, obtendo $N''$, e por fim $C_1(\phi_1)$, obtendo o referencial do corpo $B$.
% A matriz de cossenos diretores resultante é, portanto,

% \begin{equation}
%     \label{eq_sequencia_c1c2c3}
% \mathbf{C}^{B/N}
% =
% C_1(\phi_1)\,C_2(\phi_2)\,C_3(\phi_3).
% \end{equation}

% Cada rotação elementar pode ser escritas como

% \begin{equation}
% C_1(\phi_1)=
% \begin{bmatrix}
% 1 & 0 & 0 \\
% 0 & \cos\phi_1 & \sin\phi_1 \\
% 0 & -\sin\phi_1 & \cos\phi_1
% \end{bmatrix}.
% \label{eq_C1_phi1}
% \end{equation}

% \begin{equation}
% C_2(\phi_2)=
% \begin{bmatrix}
% \cos\phi_2 & 0 & -\sin\phi_2 \\
% 0 & 1 & 0 \\
% \sin\phi_2 & 0 & \cos\phi_2
% \end{bmatrix}.
% \label{eq_C2_phi2}
% \end{equation}

% \begin{equation}
% C_3(\phi_3)=
% \begin{bmatrix}
% \cos\phi_3 & \sin\phi_3 & 0 \\
% -\sin\phi_3 & \cos\phi_3 & 0 \\
% 0 & 0 & 1
% \end{bmatrix}.
% \label{eq_C3_phi3}
% \end{equation}


% As velocidades angulares associadas a cada rotação elementar são
% \begin{equation}
% \vec{\omega}_{N'/N} = \dot{\phi}_3\,\vec{n}_3.
% \label{eq_omega_Np_N}
% \end{equation}

% \begin{equation}
% \vec{\omega}_{N''/N'} = \dot{\phi}_2\,\vec{n}'_2.
% \label{eq_omega_Npp_Np}
% \end{equation}

% \begin{equation}
% \vec{\omega}_{B/N''} = \dot{\phi}_1\,\vec{b}_1.
% \label{eq_omega_B_Npp}
% \end{equation}


% O termo $\{\vec{n}_1,\vec{n}_2,\vec{n}_3\}$ é a base do referencial $N$ e $\{\vec{b}_1,\vec{b}_2,\vec{b}_3\}$ é a base do referencial $B$.

% Introduzem-se ainda os referenciais intermediários $N'$ e $N''$, associados à sequência de rotações elementares que define os ângulos de Euler.
% A velocidade angular total de $B$ em relação a $N$ pode ser obtida pela composição das velocidades angulares relativas sucessivas 

% \begin{equation}
% \vec{\omega}^{B/N}
% =
% \vec{\omega}_{B/N''}
% +
% \vec{\omega}_{N''/N'}
% +
% \vec{\omega}_{N'/N}
% =
% \dot{\phi}_1\,\vec{b}_1
% +
% \dot{\phi}_2\,\vec{n}'_2
% +
% \dot{\phi}_3\,\vec{n}_3.
% \label{eq_omega_soma_contrib}
% \end{equation}

% %%%%%%%%%%%%%%%%%%%%%%%%%%%%%%%%%%%%%%%%%%%%%

% Ao expressar a velocidade angular de $B$ em relação a $N$ na espaçonave e ao projetar a soma das contribuições elementares nos eixos de $B$, obtém-se uma relação linear entre as componentes da velocidade angular e as derivadas temporais dos ângulos de Euler.
% Essa relação pode ser escrita em forma matricial conforme apresentado por Wie~\cite{wie2008space}:

% \begin{equation}
% \begin{bmatrix}
% \omega_1 \\[2pt] \omega_2 \\[2pt] \omega_3
% \end{bmatrix}
% =
% \begin{bmatrix}
% 1 & 0 & -\sin\phi_2 \\
% 0 & \cos\phi_1 & \sin\phi_1 \cos\phi_2 \\
% 0 & -\sin\phi_1 & \cos\phi_1 \cos\phi_2
% \end{bmatrix}
% \begin{bmatrix}
% \dot{\phi}_1 \\[2pt] \dot{\phi}_2 \\[2pt] \dot{\phi}_3
% \end{bmatrix}.
% \label{eq_cinematica_euler}
% \end{equation}

% %%%%%%%%%%%%%%%%%%%%%%%%%%%%%%%%%%%%%%%%%

% A matriz $3\times 3$ em \eqref{eq_cinematica_euler} não é ortogonal, pois não está associada a uma mudança de base entre referenciais ortonormais por meio de uma rotação.
% Em particular, os vetores que sustentam essa decomposição de $\vec{\omega}$ correspondem aos eixos instantâneos das rotações elementares e, quando expressos no corpo, não formam um conjunto mutuamente ortogonal.
% Sendo assim, as direções associadas a
% $\vec{b}_1$, $\vec{a}''_2$ e $\vec{a}'_3$
% não constituem uma base ortonormal.
% Consequentemente, a matriz não representa uma rotação, mas uma transformação linear que relaciona as taxas dos ângulos de Euler às componentes de $\vec{\omega}$.


% A relação inversa,
% obtida pela inversão da Matriz~\eqref{eq_cinematica_euler},
% fornece as equações diferenciais cinemáticas para a sequência apresentada na Equação~\eqref{eq_sequencia_c1c2c3}:

% \begin{equation}
% \begin{bmatrix}
% \dot{\phi}_1 \\[2pt] \dot{\phi}_2 \\[2pt] \dot{\phi}_3
% \end{bmatrix}
% =
% \frac{1}{\cos\phi_2}
% \begin{bmatrix}
% \cos\phi_2 & \sin\phi_1 \sin\phi_2 & \cos\phi_1 \sin\phi_2 \\
% 0 & \cos\phi_1 \cos\phi_2 & -\sin\phi_1 \cos\phi_2 \\
% 0 & \sin\phi_1 & \cos\phi_1
% \end{bmatrix}
% \begin{bmatrix}
% \omega_1 \\[2pt] \omega_2 \\[2pt] \omega_3
% \end{bmatrix}.
% \label{eq_euler_inversa}
% \end{equation}

% Conforme \cite{wie2008space}, a Equação~\eqref{eq_euler_inversa} corresponde às equações cinemáticas diferenciais da sequência de Euler adotada, pois relaciona as derivadas temporais dos ângulos às componentes da velocidade angular expressas na espaçonave.
% A \gls{dcm} associada aos ângulos $(\phi_1,\phi_2,\phi_3)$ é obtida pela composição das rotações elementares definida na Equação~\eqref{eq_sequencia_c1c2c3}.
% Essa construção caracteriza uma sequência em torno dos eixos intermediários $z$--$y$--$x$ (isto é, $3$--$2$--$1$), sendo equivalente a uma sequência em torno de eixos fixos $x$--$y$--$z$, na ordem $X$--$Y$--$Z$.
% Essa equivalência permite comparar diretamente a \gls{dcm} resultante com a parametrização por ângulos de Euler empregada.

% %%%%%%%%%%%%%%%%%%%%%%%%%%%%%%%%%%%%%%%%%%%%
% A Matriz~\eqref{eq_cinematica_euler} torna-se singular quando $\phi_2=\pm 90^\circ (pitch)$, pois o mapeamento entre as derivadas dos ângulos de Euler e as componentes de $\vec{\omega}$ perde posto e a transformação inversa \eqref{eq_euler_inversa} deixa de existir.
% Nessa configuração, na sequência adotada, os eixos associados às rotações primeira e terceira tornam-se alinhados, e elimina a independência entre dois movimentos elementares e reduz o número de graus de liberdade.
% Esse fenômeno é conhecido como singularidade cinemática, ou \emph{gimbal lock}.
% A adoção de outra sequência de Euler não elimina o problema, apenas desloca a singularidade para outra combinação de ângulos.
% Essa limitação é inerente a parametrizações de atitude baseadas em três rotações sucessivas e motiva o uso de representações alternativas, como quaternions, que não apresentam singularidades de parametrização no espaço de estados.

\subsection{Ângulos de Euler}
\label{sec_angulos_euler}

Assim como a matriz de cossenos diretores, a orientação de um referencial $B$ em relação a outro $N$ pode ser descrita por ângulos de Euler.
Conforme \cite{wie2008space}, considera-se a sequência de rotações $C_1(\phi_1) \leftarrow C_2(\phi_2) \leftarrow C_3(\phi_3)$ que leva o referencial inercial $N$ até o sistema do corpo $B$.
A sequência de transformações entre os referenciais intermediários $N'$ e $N''$ é dada por:

\begin{equation}
N \xrightarrow{C_3(\phi_3)} N' \xrightarrow{C_2(\phi_2)} N'' \xrightarrow{C_1(\phi_1)} B.
\end{equation}

Isto é, parte-se do referencial inercial $N$ e aplica-se primeiro a rotação $C_3(\phi_3)$ em torno do eixo $z$, obtendo $N'$; em seguida, aplica-se $C_2(\phi_2)$ em torno do novo eixo $y'$, obtendo $N''$; e por fim, $C_1(\phi_1)$ em torno do eixo $x''$, obtendo o referencial do corpo $B$.
A matriz de cossenos diretores resultante é dada pelo produto das matrizes elementares:

\begin{equation}
\label{eq_sequencia_c1c2c3}
\mathbf{C}^{B/N}
=
C_1(\phi_1)\,C_2(\phi_2)\,C_3(\phi_3).
\end{equation}

As matrizes de rotação elementar para essa sequência são definidas como \cite{wie2008space}:

\begin{equation}
C_1(\phi_1)=
\begin{bmatrix}
1 & 0 & 0 \\
0 & \cos\phi_1 & \sin\phi_1 \\
0 & -\sin\phi_1 & \cos\phi_1
\end{bmatrix},
\label{eq_C1_phi1}
\end{equation}

\begin{equation}
C_2(\phi_2)=
\begin{bmatrix}
\cos\phi_2 & 0 & -\sin\phi_2 \\
0 & 1 & 0 \\
\sin\phi_2 & 0 & \cos\phi_2
\end{bmatrix},
\label{eq_C2_phi2}
\end{equation}

\begin{equation}
C_3(\phi_3)=
\begin{bmatrix}
\cos\phi_3 & \sin\phi_3 & 0 \\
-\sin\phi_3 & \cos\phi_3 & 0 \\
0 & 0 & 1
\end{bmatrix}.
\label{eq_C3_phi3}
\end{equation}

Os ângulos $\phi_1, \phi_2, \phi_3$ correspondem, respectivamente, aos movimentos de rolagem (\textit{roll}), arfagem (\textit{pitch}) e guinada (\textit{yaw}).
A velocidade angular total de $B$ em relação a $N$ é obtida pela soma vetorial das taxas angulares relativas em cada etapa da rotação:

\begin{equation}
\vec{\omega}^{B/N}
=
\vec{\omega}_{B/N''}
+
\vec{\omega}_{N''/N'}
+
\vec{\omega}_{N'/N}
=
\dot{\phi}_1\,\vec{b}_1
+
\dot{\phi}_2\,\vec{n}''_2
+
\dot{\phi}_3\,\vec{n}'_3.
\label{eq_omega_soma_contrib}
\end{equation}

O termo $\vec{b}_1$ é o eixo $x$ do corpo, $\vec{n}''_2$ é o eixo $y$ do referencial intermediário $N''$ e $\vec{n}'_3$ é o eixo $z$ do referencial intermediário $N'$.

Ao projetar essas contribuições nos eixos do sistema do corpo $B$ ($\{\vec{b}_1, \vec{b}_2, \vec{b}_3\}$), obtém-se a relação cinemática diferencial entre as componentes da velocidade angular $\vec{\omega}^{B/N} = [\omega_1, \omega_2, \omega_3]^T$ e as taxas dos ângulos de Euler \cite{wie2008space}:

\begin{equation}
\begin{bmatrix}
\omega_1 \\[2pt] \omega_2 \\[2pt] \omega_3
\end{bmatrix}
=
\begin{bmatrix}
1 & 0 & -\sin\phi_2 \\
0 & \cos\phi_1 & \sin\phi_1 \cos\phi_2 \\
0 & -\sin\phi_1 & \cos\phi_1 \cos\phi_2
\end{bmatrix}
\begin{bmatrix}
\dot{\phi}_1 \\[2pt] \dot{\phi}_2 \\[2pt] \dot{\phi}_3
\end{bmatrix}.
\label{eq_cinematica_euler}
\end{equation}

A relação inversa, que fornece as taxas dos ângulos de Euler em função da velocidade angular medida pelos sensores da espaçonave, é dada por:

\begin{equation}
\begin{bmatrix}
\dot{\phi}_1 \\[2pt] \dot{\phi}_2 \\[2pt] \dot{\phi}_3
\end{bmatrix}
=
\frac{1}{\cos\phi_2}
\begin{bmatrix}
\cos\phi_2 & \sin\phi_1 \sin\phi_2 & \cos\phi_1 \sin\phi_2 \\
0 & \cos\phi_1 \cos\phi_2 & -\sin\phi_1 \cos\phi_2 \\
0 & \sin\phi_1 & \cos\phi_1
\end{bmatrix}
\begin{bmatrix}
\omega_1 \\[2pt] \omega_2 \\[2pt] \omega_3
\end{bmatrix}.
\label{eq_euler_inversa}
\end{equation}

A matriz na Equação~\eqref{eq_euler_inversa} torna-se singular quando $\phi_2 = \pm 90^\circ$ (singularidade de \textit{pitch}), pois $\cos\phi_2 = 0$.
Nessa configuração, conhecida como \textit{gimbal lock}, os eixos da primeira e terceira rotações alinham-se, resultando na perda de um grau de liberdade na representação matemática da atitude.
Essa limitação motiva o uso de parametrizações globais não singulares, como os parâmetros de Euler (quatérnions), para descrever grandes manobras de atitude.